% Notes of the Montreal seminar, 1970 (cote 129, typed part): archive cote 129, site page 36 (PDF page 37).
% typescript. Draft compared: none.
% AI-transcribed by Claude Opus 5.5 (claude-opus-5-5, adaptive thinking, effort high), Message Batches API, batch mtl129_full (msgbatch_01AGnDKHwy6ATCM7ehUa6263), 2026-11-13; not yet checked by eye.
\scanpage{129/36}
\mspage{\hand{II.8}}

bi-infinitésimal) si et seulement si $G'$ et $G''$ sont infinitésimaux
(resp. étales, resp. unipotents, resp. de type multiplicatif, resp. bi-%
infinitésimaux).

On vérifie aisément l'assertion dans le cas infinitésimal
et dans le cas étale en se ramenant au cas où $S$ est le spectre d'un
corps algébriquement clos. Les autres cas s'en déduisent par dualité de
Cartier.

2.4. Exemples de p-groupes en caractéristique $p$ .

$(\mathbb{Z}/p\mathbb{Z})_S$ est étale sur $S$ et $\mu_{p|S} = (\mathbb{Z}/p\mathbb{Z})^*_S$ est donc de
type multiplicatif. Par ailleurs, $\alpha_{p|S} = (\alpha_{p|S})^*$ est bi-infinitésimal.
Parmi ces 3 groupes, $\alpha_{p|S}$ et $(\mathbb{Z}/p\mathbb{Z})_S$ sont les groupes unipotents tan-%
disque les groupes infinitésimaux sont $\alpha_{p|S}$ et $\mu_{p|S}$ .

L'intérêt de ces trois groupes réside dans le résultat sui-%
vant : si $S$ est le spectre d'un corps algébriquement clos et de carac-%
téristique $p$ , tout groupe annulé par une puissance de $p$ et fini sur
$S$ admet une suite de composition finie
\margin{$\times$}
\[ G \supset G_0 \supset G_1 \supset G_2 \ldots \supset G_{n-1} \supset G_n \hand{$= 0$} \]
où les $G_i/G_{i+1}$ sont isomorphes à l'un des groupes $\alpha_p$ , $\mu_p$ , $\mathbb{Z}/p\mathbb{Z}$ . De
plus, il résulte de 2.4. que $G$ est infinitésimal (resp. étale, resp.
unipotent, resp. de type multiplicatif, resp. bi-infinitésimal) si et
seulement si tous les $G_i/G_{i+1}$ sont isomorphes à $\alpha_p$ ou $\mu_p$ (resp.
$\mathbb{Z}/p\mathbb{Z}$ , resp. $\mathbb{Z}/p\mathbb{Z}$ ou $\alpha_p$ , resp. $\mu_p$ , resp. $\alpha_p$ ). (cf. SGA 3, XVII, 1.7.
et 4.2.1).

\margin{36}
